\section{Introduction}
\label{sec:intro}

API로 부르는 프런티어 LLM으로 로봇을 제어하는 연구가 늘고 있다~\cite{agp2026,robodojo2026,su2026astra}. 이런 모델은 과제를 나누어 계획하고 처음 보는 장면에서도 비교적 잘 버틴다. 반면 과제 데이터로 미세조정한 VLA는 환경이 바뀌면 크게 떨어진다. RoboDojo~\cite{robodojobench2026}에서 $\pi_{0.5}$의 점수는 깨끗한 장면 20.92에서 바뀐 장면 5.82로 72\% 떨어졌다. 같은 벤치마크의 공개 원자료로 우리가 계산한 GPT-6 Astra(이하 Astra)의 하락은 \tentative{11.1\%}였다(조건이 같지 않은 예비 분석). 우리는 이 차이를 다음처럼 읽는다. 짧은 관측 창으로 과제 전체를 끝-끝으로 배우는 VLA는 좁은 데이터에서 과제의 의미 대신 장면에 묶인 지름길을 배우기 쉽다~\cite{shortcut2025,liberoplus2025}. 다만 이는 좁은 데이터와 순진한 미세조정이 겹칠 때의 일이며, 데이터를 넓히거나 VLM 지식을 보존하면 새 환경을 따라간 사례도 있다~\cite{pi05_2025,kinsulate2025}.

그러나 프런티어 계획기를 제어 루프에 그대로 넣으면 두 문제가 생긴다. 첫째, 느리다. 사전 등록 탐침에서 Astra(effort low)가 카메라 세 대를 보고 답하는 데 p50 \tentative{8.9--10.0\,s}가 걸렸다. 둘째, 판단마다 로봇을 세우면 실제로 쓸 수 없고, 세우지 않으면 몇 초 전 영상으로 만든 흔들리는 답을 지금의 움직임에 어떻게 섞을지가 문제다. 기존 하네스는 상위 모델의 답을 기다리는 동기 루프이거나~\cite{harnessvla2026,su2026astra}, 가장 새 답을 그대로 쓴다~\cite{slowbrain2026}. 빠른 쪽을 학습된 작은 모델로 채우면 그 모델은 다시 부정확하다.

\paragraph{핵심 아이디어: 페이스 노트.} 랠리에서 코드라이버는 앞 구간의 페이스 노트를 미리 읽어 주고, 드라이버는 그 노트를 기준으로 조향과 제동 시점을 실시간으로 정하며, 코드라이버는 차의 위치를 확인한 뒤 다음 노트를 부른다. 둘은 서로의 일을 대신하지 않는다. PaceNotes(\cref{fig:teaser})는 이 분업을 로봇 조작에 옮긴 \emph{의도 공유 계층}이다. 상위 계획기는 행동을 낼 때마다 다음 구간의 의도(지금은 물체로 다가간다, 이번엔 잡는다)를 함께 선언한다. VLA는 상위의 답과 답 사이를 0.33\,s마다 스스로 메우되 선언된 의도를 기준으로 움직이고, 언제 쥐고 놓을지를 스스로 정한다. 상위 계획기는 VLA가 실제로 한 움직임과 결과를 보고 잡았는지·놓았는지를 검증한 뒤 다음 의도를 낸다. 이 규약은 선택 사항이 아니라 시스템의 필수 핵심이다. VLA는 과제 전체의 정책이 아니라 방향·크기·대상·단계·그리퍼를 닫힌 보기로 고르는 조이스틱이고, 행동 전문가가 고른 보기를 연속 동작으로 채운다. 느린 상위의 답은 도착 시점에 그사이 VLA가 한 일과 대조하고, 두 답의 합의와 물체와의 거리에 따른 권한을 거쳐 서서히 반영한다.

\paragraph{상위 계획기를 학습 가능한 크기로.} Astra는 한 호출에 수 초가 걸리고 호출마다 비용이 든다. 우리의 최종 목표는 계층을 유지한 채 상위 자리의 Astra를 우리가 학습할 수 있는 크기의 오픈 VLM(Qwen3.5-35B-A3B; 8B는 경향을 보는 임시 대리)으로 바꾸는 것이다. \tentative{우리는 일반 능력은 Astra보다 낮더라도 말단과 로봇·장면의 관계와 이동의 상관에 특화해 학습한 상위 모델이 더 빠르면서 우리 과제에서는 비슷하거나 더 나을 것이라고 가정한다(가설 H-L, 미확정).} 근거는 우리 측정에 있다. 영샷 오픈급 모델의 실패는 모두 대상 위치 읽기였고, Astra도 7--10\,cm 목표 어긋남은 알아채지 못했다. 둘 다 시뮬의 촘촘한 위치 감독이 겨누는 약점이다.

\paragraph{기여.} 이 원고는 가설과 승인된 설계로 쓴 완성 구조의 원고이며, 아래 주장은 사전 등록한 실험으로 검증한다. 주황색 값은 아직 확정되지 않은 값(측정 초기값, 가정값, 결과 예정 칸)이다.
\begin{enumerate}
  \item \textbf{의도 공유 계층.} 상위 계획기의 구간 의도 선언과 결과 검증, VLA의 실시간 세부 조종과 쥐기·놓기 시점 소유, 느린 답의 도착 시 대조·합의·거리 권한을 한 규약으로 묶는다(\cref{sec:protocol}). 검증할 주장: 같은 호출 예산에서 결합 계층이 VLA 단독과 동기식 상위 단독보다 성공률이 높거나, 같은 성공률에서 로봇이 멈추는 시간과 비용이 작다(\cref{sec:plan}).
  \item \textbf{계층을 위한 학습 레시피.} (a) 프런티어 계획기를 학습 가능한 크기의 상위 모델로 증류하는 선생-학생 학습, (b) 조이스틱 역할에 맞춘 VLA의 끝-끝 학습(먼 구간 반사실 분기로 결정이 실행을 조종하게 함), (c) 다양화한 시뮬과 좌표계 정보가 없는 공개 데이터의 픽셀 경로 변환(\cref{sec:training}).
  \item \textbf{초기 근거와 진단.} Astra에 맞춘 단독 조종 \tentative{7/7}, 영샷 오픈급 모델의 위치 읽기 실패, 시뮬 참값으로 학습한 8B 상위 모델의 분포 안 \tentative{4/4}와 접근 오차 \tentative{2.6\,mm}, Astra의 목표 어긋남 알아채기 공백(\tentative{0/13}), VLA 단독 \tentative{0/12}의 원인 분해(\cref{sec:evidence}).
  \item \textbf{검증 사다리와 평가 프로토콜.} 결합 폐루프까지의 무료 단계별 관문, 같은 장면 짝(standard, random)의 상대 하락(RD)으로 위층과 아래층의 기여를 가르는 평가, 실행 전에 고정한 판정 기준과 그 코드를 대조하는 독립 감사(\cref{sec:plan}).
\end{enumerate}
